%============================================================
%  Bd. 41 - Endliche Gruppen
%============================================================
\documentclass[main.tex]{subfiles}
\ifSubfilesClassLoaded{\BandSetupStandalone{B41}}{}
\title{Bd. 41 - Endliche Gruppen}
\author{Martin Kunze}
\date{}
\setcounter{file}{41}
\hypersetup{pdftitle={Bd. 41 - Endliche Gruppen},pdfauthor={Martin Kunze}}
\begin{document}
\title{Bd. 41 - Endliche Gruppen}
\author{Martin Kunze}
\date{}
\hypersetup{pageanchor=false}
\maketitle
\hypersetup{pageanchor=true}
\setcounter{tocdepth}{3}
\tableofcontents
\setcounter{file}{41}
\renewcommand{\FormulaBandID}{41}

\chapter{Einleitung}

Dieser Band schließt an die Gruppen aus Band~40 an. Sein erster Schwerpunkt
ist eine Anwendung endlicher Gruppen auf das Isomorphieproblem aus Band~37:
Eine endliche Halbgruppe ist durch ihre Potenzhalbgruppe bestimmt, sobald
ihre Produktmenge eine Gruppe ist.

Die Beweisidee hat zwei Teile. Zunächst wird die Gruppe im Produktquadrat
erkannt. Danach wird gezählt, wie viele Elemente der Ausgangshalbgruppe auf
jedes Gruppenelement zurückgezogen werden. Die Potenzhalbgruppe kennt diese
Anzahlen als Größen von Lösungsmengen. Die Endlichkeit erlaubt, von den
Potenzmengen wieder zu den Grundfasern zurückzukehren.

Die benötigten allgemeinen Werkzeuge stehen dort, wo ihre Voraussetzungen
erstmals verfügbar sind: punktierte Bijektionen in Band~08, das Verkleben
von Faserbijektionen in Band~09, endliche Potenzmengenkardinalitäten in
Band~20, Produktquadrate und Inflationen in Band~28, monoidale Quadrate in
Band~38 sowie der Transport von Gruppeneinermengen in Band~40.
Der vorliegende Band verbindet diese Bausteine; er führt keine zusätzliche
Annahme zur Potenzhalbgruppenvermutung ein.

\chapter{Endliche Gruppen}

\section{Begriff und Einordnung}

\FormulaDefDeltaK[Begriff der endlichen Gruppe]{%
  \mathsf{FinGrp}(G,\star,e)%
}{FiniteGroupDef}{%
  \DeltaRow{Mengen}{G\dsep e}
  \DeltaRow{Funktionen}{\star}
  \DeltaRow{\textbf{Axiome}}{}
  \DeltaPrem{Gruppe}{\GroupAlg{G}{\star}{e}}
  \DeltaPrem{Endlichkeit}{\Finite{G}}
  \DeltaRow{\textbf{Neues Symbol}}{}
  \DeltaPrem{Endliche Gruppe}{\mathsf{FinGrp}(G,\star,e)}
}

\FormulaAxiomDeltaK[Zugriff auf die Gruppenstruktur]{%
  \mathsf{FinGrp}(G,\star,e)\vdash\GroupAlg{G}{\star}{e}%
}{FiniteGroupBaseAxiom}{%
  \DeltaRow{Mengen}{G\dsep e}
  \DeltaRow{Funktionen}{\star}
}

\FormulaAxiomDeltaK[Zugriff auf die Endlichkeit]{%
  \mathsf{FinGrp}(G,\star,e)\vdash\Finite{G}%
}{FiniteGroupFinitenessAxiom}{%
  \DeltaRow{Mengen}{G\dsep e}
  \DeltaRow{Funktionen}{\star}
}

\par\noindent\begin{minipage}{\linewidth}
\FormulaThmDeltaK[Eine endliche Gruppe ist eine endliche Halbgruppe]{%
  \mathsf{FinGrp}(G,\star,e)\vdash\FiniteSemigroupAlg{G}{\star}%
}{FiniteGroupIsFiniteSemigroup}{%
  \DeltaRow{Mengen}{G\dsep e}
  \DeltaRow{Funktionen}{\star}
}
Die Gruppenstruktur liefert die Halbgruppenstruktur; die Endlichkeit
des Trägers bleibt dabei erhalten.
\end{minipage}\par\nopagebreak[4]
\begin{tabproofpaired}
  \proofsteppaired[1]{\mathsf{FinGrp}(G,\star,e)}{\rA}
  \proofsteppaired[1]{\GroupAlg{G}{\star}{e}}{%
    \FormulaRefAuto{FiniteGroupBaseAxiom}[ax]{1}}
  \proofsteppaired[1]{\Finite{G}}{%
    \FormulaRefAuto{FiniteGroupFinitenessAxiom}[ax]{1}}
  \proofsteppaired[1]{\SemigroupAlg{G}{\star}}{%
    \FormulaRefAuto{GroupBaseSemigroup}[thm]{2}}
  \proofsteppaired[1]{\FiniteSemigroupAlg{G}{\star}}{%
    \FormulaRefAuto{FiniteSemigroupDef}[def]{4,3}}
\end{tabproofpaired}

\chapter{Rekonstruktion bei einem Gruppenquadrat}
\section{Notation und der erkennbare Gruppenkern}

Wir verwenden das Produktquadrat aus Band~28:
\[
  S^2=S\star_{\mathcal P}S=\{x\star y\mid x,y\in S\}.
\]
Es ist vom kartesischen Produkt \(S\times S\) zu unterscheiden.
Auf Unterhalbgruppen bezeichnet dieselbe Operationsschreibweise stets die
entsprechende Einschränkung. Für die folgenden Aussagen setzen wir
\[
  \begin{gathered}
    G=S^2,\qquad H=T^2,\qquad F_0=F\restriction_{Q_S^2},\\[-2pt]
    Q_S=\PowNE S,\qquad Q_T=\PowNE T,\\[-2pt]
    Q_G=\PowNE G,\qquad Q_H=\PowNE H.
  \end{gathered}
\]
Die Operationen auf \(S\) und \(T\) heißen \(\star\) und \(\diamond\),
die Mengenprodukte \(\star_{\mathcal P}\) und \(\diamond_{\mathcal P}\).
Die hier auftretenden Quadrate werden auf nichtleeren Trägern gebildet.

\paragraph{Beweisidee.}
Ein Isomorphismus erhält das Produktquadrat. Dieses ist auf der Quellseite
die ganze Potenzhalbgruppe von \(G\). Seine Monoidstruktur erkennt auf der
Zielseite zunächst \(H\) als Monoid. Dann greifen die bereits bewiesene
Gruppenstarrheit und der Transport der Gruppeneinermengen.

\par\noindent\begin{minipage}{\linewidth}
\FormulaThmDeltaK[Ein Gruppenquadrat wird mitsamt seinen Einermengen rekonstruiert]{%
  \begin{aligned}[t]
    &\SemigroupAlg{S}{\star}\dsep\SemigroupAlg{T}{\diamond}
      \dsep S\neq\varnothing\dsep T\neq\varnothing\dsep{}\\[-2pt]
    &\GroupAlg{G}{\star}{e}\dsep
      \SemigroupIso{F}{Q_S,\star_{\mathcal P}}{Q_T,\diamond_{\mathcal P}}
      \vdash{}\\[-2pt]
    &\exists u\in H\,\exists\varphi\,\left(
      \begin{aligned}[t]
        &\GroupAlg{H}{\diamond}{u}\land
          \GroupIso{\varphi}{G,\star,e}{H,\diamond,u}\land{}\\[-2pt]
        &\forall g\in G\,(F(\{g\})=\{\varphi(g)\})
      \end{aligned}\right)
  \end{aligned}
}{PowerGroupSquareCoreReconstruction}{%
  \DeltaRow{Mengen}{S\dsep T\dsep G=S^2\dsep H=T^2\dsep e\dsep u\dsep v\dsep g}
  \DeltaRow{Potenzträger}{Q_S=\PowNE S\dsep Q_T=\PowNE T}
  \DeltaRow{Funktionen}{\star\dsep\diamond\dsep F\dsep F_0\dsep\varphi\dsep h}
}
\end{minipage}\par
\begin{tabproofsplitwide}
  \proofpartpairedindR[Das Zielquadrat ist ein Monoid]{%
    \begin{aligned}[t]
      &\SemigroupAlg{S}{\star}\dsep\SemigroupAlg{T}{\diamond}
        \dsep S\neq\varnothing\dsep T\neq\varnothing\dsep{}\\[-2pt]
      &\GroupAlg{G}{\star}{e}\dsep
        \SemigroupIso{F}{Q_S,\star_{\mathcal P}}{Q_T,\diamond_{\mathcal P}}
        \vdash{}\\[-2pt]
      &\exists u\in H\,\bigl(\MonoidAlg{H}{\diamond}{u}\land
        \SemigroupIso{F_0}{\PowNE G,\star_{\mathcal P}}{\PowNE H,\diamond_{\mathcal P}}\bigr)
    \end{aligned}
  }{%
    \SemigroupAlg{S}{\star}\dsep\SemigroupAlg{T}{\diamond}
      \dsep S\neq\varnothing\dsep T\neq\varnothing\dsep
      \GroupAlg{G}{\star}{e}\dsep
      \SemigroupIso{F}{Q_S,\star_{\mathcal P}}{Q_T,\diamond_{\mathcal P}}
      \vdash\exists u\in H\,\bigl(\MonoidAlg{H}{\diamond}{u}\land
      \SemigroupIso{F_0}{\PowNE G,\star_{\mathcal P}}{\PowNE H,\diamond_{\mathcal P}}\bigr)
  }[PowerGroupSquareCoreMonoidTransport]
    \proofsteppaired[1]{\SemigroupAlg{S}{\star}\land\SemigroupAlg{T}{\diamond}}{\rA}
    \proofsteppaired[2]{S\neq\varnothing\land T\neq\varnothing}{\rA}
    \proofsteppaired[3]{\GroupAlg{G}{\star}{e}}{\rA}
    \proofsteppaired[4]{\SemigroupIso{F}{Q_S,\star_{\mathcal P}}{Q_T,\diamond_{\mathcal P}}}{\rA}
    \proofsteppaired[1]{\begin{aligned}[t]
      &\SemigroupAlg{Q_S}{\star_{\mathcal P}}\\[-2pt]
      &\land\SemigroupAlg{Q_T}{\diamond_{\mathcal P}}
    \end{aligned}}{\FormulaRefAuto{NonemptyPowerSemigroup}[thm]{1}}
    \proofsteppaired[2]{Q_S\neq\varnothing\land Q_T\neq\varnothing}{%
      \FormulaRefAuto{NonemptyPowerCarrierNonempty}[thm]{2}}
    \proofsteppaired[1,2,4]{\SemigroupIso{F_0}{Q_S^2,\star_{\mathcal P}}{Q_T^2,\diamond_{\mathcal P}}}{%
      \FormulaRefAuto{SemigroupIsoSquareRestriction}[thm]{4,\allowbreak 6}}
    \proofsteppaired[3]{\MonoidAlg{G}{\star}{e}}{\FormulaRefAuto{GroupBaseMonoidAxiom}[ax]{3}}
    \proofsteppaired[1,2,3]{\begin{aligned}[t]
      &Q_S^2=\PowNE G\\[-2pt]
      &\land\MonoidAlg{Q_S^2}{\star_{\mathcal P}}{\{e\}}
    \end{aligned}}{%
      \FormulaRefAuto{PowerSemigroupSquareMonoidForward}[thm]{1,\allowbreak 2,\allowbreak 8}}
    \proofsteppaired[1,2,3,4]{\MonoidAlg{Q_T^2}{\diamond_{\mathcal P}}{F_0(\{e\})}}{%
      \FormulaRefAuto{SemigroupIsoMonoidTransport}[thm]{\rAEb{9},\allowbreak 7}}
    \proofsteppaired[1,2,3,4]{\begin{aligned}[t]
      &\exists u\in H\,\bigl(F_0(\{e\})=\{u\}\\[-2pt]
      &\quad\land\MonoidAlg{H}{\diamond}{u}\bigr)
    \end{aligned}}{%
      \FormulaRefAuto{PowerSemigroupSquareMonoidBackward}[thm]{1,\allowbreak 2,\allowbreak 10}}
    \proofsteppaired[12]{\begin{aligned}[t]
      &u\in H\land F_0(\{e\})=\{u\}\\[-2pt]
      &\land\MonoidAlg{H}{\diamond}{u}
    \end{aligned}}{\rA}
    \proofsteppaired[1,2,12]{Q_T^2=\PowNE H}{%
      \rAEa{\FormulaRefAuto{PowerSemigroupSquareMonoidForward}[thm]{1,\allowbreak 2,\allowbreak 12}}}
    \proofsteppaired[1,2,3,4,12]{%
      \SemigroupIso{F_0}{Q_G,\star_{\mathcal P}}{Q_H,\diamond_{\mathcal P}}}{\rIE{7,\allowbreak 9,\allowbreak 13}}
    \proofsteppaired[1,2,3,4,12]{\begin{aligned}[t]
      &\exists u\in H\,\bigl(\MonoidAlg{H}{\diamond}{u}\\[-2pt]
      &\quad\land\SemigroupIso{F_0}{Q_G,\star_{\mathcal P}}{Q_H,\diamond_{\mathcal P}}\bigr)
    \end{aligned}}{%
      \rEI{\rAI{\rAEa{12},\allowbreak \rAI{\rAEb{\rAEb{12}},\allowbreak 14}}}}
    \proofsteppaired[1,2,3,4]{\begin{aligned}[t]
      &\exists u\in H\,\bigl(\MonoidAlg{H}{\diamond}{u}\\[-2pt]
      &\quad\land\SemigroupIso{F_0}{Q_G,\star_{\mathcal P}}{Q_H,\diamond_{\mathcal P}}\bigr)
    \end{aligned}}{%
      \rEE{11,\allowbreak 12,\allowbreak 15}}
  \closeproofpartpairedindR

  \clearpage
  \proofpartpaired{Gruppenstarrheit und Einermengentransport}
    \proofsteppaired[1]{\SemigroupAlg{S}{\star}\land\SemigroupAlg{T}{\diamond}}{\rA}
    \proofsteppaired[2]{S\neq\varnothing\land T\neq\varnothing}{\rA}
    \proofsteppaired[3]{\GroupAlg{G}{\star}{e}}{\rA}
    \proofsteppaired[4]{\SemigroupIso{F}{Q_S,\star_{\mathcal P}}{Q_T,\diamond_{\mathcal P}}}{\rA}
    \proofsteppaired[1,2,3,4]{\begin{aligned}[t]
      &\exists u\in H\,\bigl(\MonoidAlg{H}{\diamond}{u}\\[-2pt]
      &\quad\land\SemigroupIso{F_0}{Q_G,\star_{\mathcal P}}{Q_H,\diamond_{\mathcal P}}\bigr)
    \end{aligned}}{%
      \FormulaRefAuto{PowerGroupSquareCoreMonoidTransport}[thm]{1,\allowbreak 2,\allowbreak 3,\allowbreak 4}}
    \proofsteppaired[6]{\begin{aligned}[t]
      &u\in H\land\MonoidAlg{H}{\diamond}{u}\\[-2pt]
      &\land\SemigroupIso{F_0}{Q_G,\star_{\mathcal P}}{Q_H,\diamond_{\mathcal P}}
    \end{aligned}}{\rA}
    \proofsteppaired[6]{\SemigroupAlg{H}{\diamond}}{\FormulaRefAuto{MonoidBaseSemigroupAxiom}[ax]{6}}
    \proofsteppaired[3,6]{\begin{aligned}[t]
      &\exists v\in H\,\exists h\,\bigl(\GroupAlg{H}{\diamond}{v}\\[-2pt]
      &\quad\land\GroupIso{h}{G,\star,e}{H,\diamond,v}\bigr)
    \end{aligned}}{%
      \FormulaRefAuto{GroupPowerSemigroupRigidity}[thm]{3,\allowbreak 7,\allowbreak 6}}
    \proofsteppaired[9]{\begin{aligned}[t]
      &v\in H\land\exists h\,\bigl(\GroupAlg{H}{\diamond}{v}\\[-2pt]
      &\quad\land\GroupIso{h}{G,\star,e}{H,\diamond,v}\bigr)
    \end{aligned}}{\rA}
    \proofsteppaired[9]{\begin{aligned}[t]
      &\exists h\,\bigl(\GroupAlg{H}{\diamond}{v}\\[-2pt]
      &\quad\land\GroupIso{h}{G,\star,e}{H,\diamond,v}\bigr)
    \end{aligned}}{\rAEb{9}}
    \proofsteppaired[11]{\begin{aligned}[t]
      &\GroupAlg{H}{\diamond}{v}\\[-2pt]
      &\land\GroupIso{h}{G,\star,e}{H,\diamond,v}
    \end{aligned}}{\rA}
    \proofsteppaired[11]{\GroupAlg{H}{\diamond}{v}}{\rAEa{11}}
    \proofsteppaired[9]{\GroupAlg{H}{\diamond}{v}}{\rEE{10,\allowbreak 11,\allowbreak 12}}
    \proofsteppaired[3,6,9]{\begin{aligned}[t]
      &\exists\varphi\,\bigl(\GroupIso{\varphi}{G,\star,e}{H,\diamond,v}\\[-2pt]
      &\quad\land\forall g\in G\,(F_0(\{g\})=\{\varphi(g)\})\bigr)
    \end{aligned}}{%
      \FormulaRefAuto{PowerGroupsSingletonTransport}[thm]{3,\allowbreak 13,\allowbreak 6}}
    \proofsteppaired[15]{\begin{aligned}[t]
      &\GroupIso{\varphi}{G,\star,e}{H,\diamond,v}\\[-2pt]
      &\land\forall g\in G\,(F_0(\{g\})=\{\varphi(g)\})
    \end{aligned}}{\rA}
    \proofsteppaired[16]{g\in G}{\rA}
    \proofsteppaired[16]{\{g\}\in\PowNE G}{\FormulaRefAuto{MonoidPowerSingletonCarrier}[thm]{16}}
    \proofsteppaired[1,2,3]{Q_S^2=\PowNE G}{%
      \rAEa{\FormulaRefAuto{PowerSemigroupSquareMonoidForward}[thm]{1,\allowbreak 2,\allowbreak 
        \FormulaRefAuto{GroupBaseMonoidAxiom}[ax]{3}}}}
    \proofsteppaired[2]{Q_S\neq\varnothing}{\FormulaRefAuto{NonemptyPowerCarrierNonempty}[thm]{2}}
    \proofsteppaired[1,2]{Q_S^2\subseteq Q_S}{%
      \FormulaRefAuto{SemigroupSquareSubset}[thm]{\FormulaRefAuto{NonemptyPowerSemigroup}[thm]{1},\allowbreak 19}}
    \proofsteppaired[4]{F\colon Q_S\to Q_T}{%
      \FormulaRefAuto{Bijektive Funktion}[def]{\FormulaRefAuto{SemigroupIsoBijectiveAxiom}[ax]{4}}}
    \proofsteppaired[1,2,3,4,16]{F_0(\{g\})=F(\{g\})}{%
      \FormulaRefAuto{C\subseteq A\dsep x\in C\vdash F\restriction_C(x)=F(x)}[thm]{21,\allowbreak 20,\allowbreak \rIE{18,\allowbreak 17}}}
    \proofsteppaired[15,16]{F_0(\{g\})=\{\varphi(g)\}}{\rRE{\rUE{\rAEb{15}},\allowbreak 16}}
    \proofsteppaired[1,2,3,4,15,16]{F(\{g\})=\{\varphi(g)\}}{\rIE{22,\allowbreak 23}}
    \proofsteppaired[1,2,3,4,15]{\forall g\in G\,(F(\{g\})=\{\varphi(g)\})}{\rUI{\rRI{16,\allowbreak 24}}}
    \proofsteppaired[1,2,3,4,9,15]{\begin{aligned}[t]
      &\exists u\in H\,\exists\varphi\,\bigl(\GroupAlg{H}{\diamond}{u}\\[-2pt]
      &\quad\land\GroupIso{\varphi}{G,\star,e}{H,\diamond,u}\\[-2pt]
      &\quad\land\forall g\in G\,(F(\{g\})=\{\varphi(g)\})\bigr)
    \end{aligned}}{%
      \rEI{\rAI{\rAEa{9},\allowbreak \rEI{\rAI{13,\allowbreak \rAI{\rAEa{15},\allowbreak 25}}}}}}
    \proofsteppaired[1,2,3,4,6,9]{\begin{aligned}[t]
      &\exists u\in H\,\exists\varphi\,\bigl(\GroupAlg{H}{\diamond}{u}\\[-2pt]
      &\quad\land\GroupIso{\varphi}{G,\star,e}{H,\diamond,u}\\[-2pt]
      &\quad\land\forall g\in G\,(F(\{g\})=\{\varphi(g)\})\bigr)
    \end{aligned}}{\rEE{14,\allowbreak 15,\allowbreak 26}}
    \proofsteppaired[1,2,3,4,6]{\begin{aligned}[t]
      &\exists u\in H\,\exists\varphi\,\bigl(\GroupAlg{H}{\diamond}{u}\\[-2pt]
      &\quad\land\GroupIso{\varphi}{G,\star,e}{H,\diamond,u}\\[-2pt]
      &\quad\land\forall g\in G\,(F(\{g\})=\{\varphi(g)\})\bigr)
    \end{aligned}}{\rEE{8,\allowbreak 9,\allowbreak 27}}
    \proofsteppaired[1,2,3,4]{\begin{aligned}[t]
      &\exists u\in H\,\exists\varphi\,\bigl(\GroupAlg{H}{\diamond}{u}\\[-2pt]
      &\quad\land\GroupIso{\varphi}{G,\star,e}{H,\diamond,u}\\[-2pt]
      &\quad\land\forall g\in G\,(F(\{g\})=\{\varphi(g)\})\bigr)
    \end{aligned}}{\rEE{5,\allowbreak 6,\allowbreak 28}}
  \closeproofpartpaired
\end{tabproofsplitwide}

\section{Von Lösungsmengen zu endlichen Fasern}

Für monoidale Quadrate mit Identitäten \(e\) und \(u\) bezeichnen
\(r=r_e\) und \(s=r_u\) die Funktionen aus
\FormulaRefAuto{MonoidSquareInflationDef}[def]. Wir schreiben
\[
  r(x)=e\star x,\qquad s(y)=u\diamond y,
  \qquad C_g=r^{-1}[\{g\}],\quad D_h=s^{-1}[\{h\}].
\]

\paragraph{Beweisidee.}
Die Gleichung \(\{e\}\star_{\mathcal P}X=\{g\}\) sagt genau,
dass \(X\) eine nichtleere Teilmenge der Faser \(C_g\) ist.
Ein Isomorphismus transportiert ihre Lösungsmenge. Die endliche
Potenzmengenzählung liefert dann die gewünschte Bijektion der Grundfasern.

\par\noindent\begin{minipage}{\linewidth}
\FormulaThmDeltaK[Endliche Inflationsfasern werden gleichmächtig]{%
  \begin{aligned}[t]
    &\FiniteSemigroupAlg{S}{\star}\dsep\FiniteSemigroupAlg{T}{\diamond}
      \dsep S\neq\varnothing\dsep T\neq\varnothing\dsep{}\\[-2pt]
    &\GroupAlg{G}{\star}{e}\dsep\GroupAlg{H}{\diamond}{u}\dsep
      \GroupIso{\varphi}{G,\star,e}{H,\diamond,u}\dsep{}\\[-2pt]
    &\SemigroupIso{F}{Q_S,\star_{\mathcal P}}{Q_T,\diamond_{\mathcal P}}\dsep
      \forall g\in G\,(F(\{g\})=\{\varphi(g)\})\vdash{}\\[-2pt]
    &\forall g\in G\,\exists f_g\colon C_g\bij D_{\varphi(g)}
  \end{aligned}
}{FiniteGroupSquareFiberBijections}{%
  \DeltaRow{Mengen}{S\dsep T\dsep G=S^2\dsep H=T^2\dsep e\dsep u\dsep g\dsep h\dsep x\dsep y}
  \DeltaRow{Potenzträger}{Q_S=\PowNE S\dsep Q_T=\PowNE T}
  \DeltaRow{Fasern}{C_g=r^{-1}[\{g\}]\dsep D_h=s^{-1}[\{h\}]}
  \DeltaRow{Funktionen}{\star\dsep\diamond\dsep F\dsep\varphi\dsep r\dsep s\dsep f_g}
}
\end{minipage}\par
\begin{tabproofpaired}
  \proofsteppaired[1]{\FiniteSemigroupAlg{S}{\star}\land\FiniteSemigroupAlg{T}{\diamond}}{\rA}
  \proofsteppaired[2]{S\neq\varnothing\land T\neq\varnothing}{\rA}
  \proofsteppaired[3]{\GroupAlg{G}{\star}{e}\land\GroupAlg{H}{\diamond}{u}}{\rA}
  \proofsteppaired[4]{\GroupIso{\varphi}{G,\star,e}{H,\diamond,u}}{\rA}
  \proofsteppaired[5]{\SemigroupIso{F}{Q_S,\star_{\mathcal P}}{Q_T,\diamond_{\mathcal P}}}{\rA}
  \proofsteppaired[6]{\forall g\in G\,(F(\{g\})=\{\varphi(g)\})}{\rA}
  \proofsteppaired[7]{g\in G}{\rA}
  \proofsteppaired[1]{\SemigroupAlg{S}{\star}\land\SemigroupAlg{T}{\diamond}}{%
    \FormulaRefAuto{FiniteSemigroupBaseAxiom}[ax]{1}}
  \proofsteppaired[3]{\MonoidAlg{G}{\star}{e}\land\MonoidAlg{H}{\diamond}{u}}{%
    \FormulaRefAuto{GroupBaseMonoidAxiom}[ax]{3}}
  \proofsteppaired[4]{\varphi\colon G\to H}{%
    \FormulaRefAuto{GroupHomFunctionAxiom}[ax]{\FormulaRefAuto{GroupIsoDef}[def]{4}}}
  \proofsteppaired[4,7]{\varphi(g)\in H}{%
    \FormulaRefAuto{F\colon A\to B\dsep x\in A\vdash F(x)\in B}[thm]{10,\allowbreak 7}}
  \proofsteppaired[3]{e\in G\land u\in H}{\FormulaRefAuto{MonoidIdentityCarrierAxiom}[ax]{9}}
  \proofsteppaired[4]{\varphi(e)=u}{%
    \FormulaRefAuto{GroupHomIdentityPreservation}[thm]{\FormulaRefAuto{GroupIsoDef}[def]{4}}}
  \proofsteppaired[3,4,6]{F(\{e\})=\{u\}}{\rIE{13,\allowbreak \rRE{\rUE{6},\allowbreak \rAEa{12}}}}
  \proofsteppaired[6,7]{F(\{g\})=\{\varphi(g)\}}{\rRE{\rUE{6},\allowbreak 7}}
  \proofsteppaired[1,2]{G\subseteq S\land H\subseteq T}{%
    \FormulaRefAuto{SemigroupSquareSubset}[thm]{8,\allowbreak 2}}
  \proofsteppaired[1,2,3,4,7]{\begin{aligned}[t]
      &e\in S\land g\in S\\[-2pt]
      &\land u\in T\land\varphi(g)\in T
    \end{aligned}}{%
    \FormulaRefAuto{A\subseteq B\dsep x\in A\vdash x\in B}[thm]{16,\allowbreak 12,\allowbreak 7,\allowbreak 11}}
  \proofsteppaired[1,2,3,4,7]{\begin{aligned}[t]
      &\{e\},\{g\}\in Q_S\\[-2pt]
      &\land\{u\},\{\varphi(g)\}\in Q_T
    \end{aligned}}{%
    \FormulaRefAuto{MonoidPowerSingletonCarrier}[thm]{17}}
  \proofsteppaired[1,2,3]{r\colon S\to G\land s\colon T\to H}{%
    \FormulaRefAuto{MonoidSquareInflationDef}[def]{8,\allowbreak 2,\allowbreak 9}}
  \proofsteppaired[1,2,3,4,7]{%
    \begin{aligned}[t]
      &C_g=\{x\in S\mid r(x)=g\},\\[-2pt]
      &D_{\varphi(g)}=\{y\in T\mid s(y)=\varphi(g)\}
    \end{aligned}}{%
    \FormulaRefAuto{F\colon A\to B\dsep y\in B\vdash
      F^{-1}[\{y\}]=\{\,x\in A\mid F(x)=y\,\}}[thm]{19,\allowbreak 7,\allowbreak 11}}
  \proofsteppaired[1,2,3,7]{L_{Q_S}(\{e\},\{g\})=\PowNE{C_g}}{%
    \FormulaRefAuto{MonoidSquarePowerFiber}[thm]{8,\allowbreak 2,\allowbreak 9,\allowbreak 7};\allowbreak \
    \FormulaRefAuto{SemigroupLeftEquationFiberDef}[def]{\allowbreak\FormulaRefAuto{NonemptyPowerSemigroup}[thm]{8},\allowbreak 18};\allowbreak \ \rIE{20}}
  \proofsteppaired[1,2,3,4,7]{\begin{aligned}[t]
      &L_{Q_T}(\{u\},\{\varphi(g)\})\\[-2pt]
      &{}=\PowNE{D_{\varphi(g)}}
    \end{aligned}}{%
    \FormulaRefAuto{MonoidSquarePowerFiber}[thm]{8,\allowbreak 2,\allowbreak 9,\allowbreak 11};\allowbreak \
    \FormulaRefAuto{SemigroupLeftEquationFiberDef}[def]{\allowbreak\FormulaRefAuto{NonemptyPowerSemigroup}[thm]{8},\allowbreak 18};\allowbreak \ \rIE{20}}
  \proofsteppaired[1,2,3,4,5,7]{\begin{aligned}[t]
      &F\restriction_{L_{Q_S}(\{e\},\{g\})}\colon\\[-2pt]
      &\quad L_{Q_S}(\{e\},\{g\})\\[-2pt]
      &\quad\bij L_{Q_T}(F(\{e\}),F(\{g\}))
    \end{aligned}}{%
    \FormulaRefAuto{SemigroupLeftEquationFiberTransport}[thm]{5,\allowbreak 18}}
  \proofsteppaired[1,2,3,4,5,6,7]{\begin{aligned}[t]
      &F\restriction_{\PowNE{C_g}}\colon\\[-2pt]
      &\quad\PowNE{C_g}\bij\PowNE{D_{\varphi(g)}}
    \end{aligned}}{%
    \rIE{23,\allowbreak 14,\allowbreak 15,\allowbreak 21,\allowbreak 22}}
  \proofsteppaired[1,2,3,4,5,6,7]{\PowNE{C_g}\EqCard\PowNE{D_{\varphi(g)}}}{%
    \FormulaRefAuto{A\EqCard B\coloneqq\exists F\colon A\bij B}[def]{\rEI{24}}}
  \proofsteppaired[1,2,3,4,7]{C_g\subseteq S\land D_{\varphi(g)}\subseteq T}{%
    \FormulaRefAuto{\{x\in A\mid P(x)\}\subseteq A}[thm];\allowbreak \ \rIE{20}}
  \proofsteppaired[1,2,3,4,7]{\Finite{C_g}\land\Finite{D_{\varphi(g)}}}{%
    \FormulaRefAuto{FiniteSubset}[thm]{\FormulaRefAuto{FiniteSemigroupFinitenessAxiom}[ax]{1},\allowbreak 26}}
  \proofsteppaired[1,2,3,4,5,6,7]{\exists f_g\colon C_g\bij D_{\varphi(g)}}{%
    \FormulaRefAuto{FiniteNonemptyPowerSetReflectsBijection}[thm]{27,\allowbreak 25}}
  \proofsteppaired[1,2,3,4,5,6]{\forall g\in G\,\exists f_g\colon C_g\bij D_{\varphi(g)}}{%
    \rUI{\rRI{7,\allowbreak 28}}}
\end{tabproofpaired}

\section{Der Rekonstruktionssatz}

\paragraph{Beweisidee.}
Der erste Hilfssatz liefert den Isomorphismus der Gruppenkerne und seinen
Einermengentransport. Der zweite liefert passende Faserbijektionen.
Die Inflationsrekonstruktion aus Band~28 verbindet beide Aussagen.

\par\noindent\begin{minipage}{\linewidth}
\FormulaThmDeltaK[Endliche Halbgruppen mit Gruppenquadrat sind global bestimmt]{%
  \begin{aligned}[t]
    &\FiniteSemigroupAlg{S}{\star}\dsep\FiniteSemigroupAlg{T}{\diamond}
      \dsep S\neq\varnothing\dsep T\neq\varnothing\dsep{}\\[-2pt]
    &\GroupAlg{S^2}{\star}{e}\dsep
      \SemigroupIso{F}{\PowNE S,\star_{\mathcal P}}{\PowNE T,\diamond_{\mathcal P}}
      \vdash{}\\[-2pt]
    &\exists f\,\SemigroupIso{f}{S,\star}{T,\diamond}
  \end{aligned}
}{FiniteGroupSquarePowerSemigroupReconstruction}{%
  \DeltaRow{Mengen}{S\dsep T\dsep G=S^2\dsep H=T^2\dsep e\dsep u}
  \DeltaRow{Potenzträger}{Q_S=\PowNE S\dsep Q_T=\PowNE T}
  \DeltaRow{Funktionen}{\star\dsep\diamond\dsep F\dsep\varphi\dsep r\dsep s\dsep f}
}
\end{minipage}\par
\begin{tabproofpaired}
  \proofsteppaired[1]{\FiniteSemigroupAlg{S}{\star}\land\FiniteSemigroupAlg{T}{\diamond}}{\rA}
  \proofsteppaired[2]{S\neq\varnothing\land T\neq\varnothing}{\rA}
  \proofsteppaired[3]{\GroupAlg{G}{\star}{e}}{\rA}
  \proofsteppaired[4]{\SemigroupIso{F}{Q_S,\star_{\mathcal P}}{Q_T,\diamond_{\mathcal P}}}{\rA}
  \proofsteppaired[1]{\SemigroupAlg{S}{\star}\land\SemigroupAlg{T}{\diamond}}{%
    \FormulaRefAuto{FiniteSemigroupBaseAxiom}[ax]{1}}
  \proofsteppaired[1,2,3,4]{\begin{aligned}[t]
      &\exists u\in H\,\exists\varphi\,\bigl(\GroupAlg{H}{\diamond}{u}\\[-2pt]
      &\quad\land\GroupIso{\varphi}{G,\star,e}{H,\diamond,u}\\[-2pt]
      &\quad\land\forall g\in G\,(F(\{g\})=\{\varphi(g)\})\bigr)
    \end{aligned}}{\FormulaRefAuto{PowerGroupSquareCoreReconstruction}[thm]{5,\allowbreak 2,\allowbreak 3,\allowbreak 4}}
  \proofsteppaired[7]{\begin{aligned}[t]
      &u\in H\land\exists\varphi\,\bigl(\GroupAlg{H}{\diamond}{u}\\[-2pt]
      &\quad\land\GroupIso{\varphi}{G,\star,e}{H,\diamond,u}\\[-2pt]
      &\quad\land\forall g\in G\,(F(\{g\})=\{\varphi(g)\})\bigr)
    \end{aligned}}{\rA}
  \proofsteppaired[7]{\begin{aligned}[t]
      &\exists\varphi\,\bigl(\GroupAlg{H}{\diamond}{u}\\[-2pt]
      &\quad\land\GroupIso{\varphi}{G,\star,e}{H,\diamond,u}\\[-2pt]
      &\quad\land\forall g\in G\,(F(\{g\})=\{\varphi(g)\})\bigr)
    \end{aligned}}{\rAEb{7}}
  \proofsteppaired[9]{\begin{aligned}[t]
      &\GroupAlg{H}{\diamond}{u}\\[-2pt]
      &\land\GroupIso{\varphi}{G,\star,e}{H,\diamond,u}\\[-2pt]
      &\land\forall g\in G\,(F(\{g\})=\{\varphi(g)\})
    \end{aligned}}{\rA}
  \proofsteppaired[3,9]{\MonoidAlg{G}{\star}{e}\land\MonoidAlg{H}{\diamond}{u}}{%
    \FormulaRefAuto{GroupBaseMonoidAxiom}[ax]{3,\allowbreak 9}}
  \proofsteppaired[1,2,3,9]{\begin{aligned}[t]
      &\mathsf{Infl}_{r}((G,\star),(S,\star))\\[-2pt]
      &\land\mathsf{Infl}_{s}((H,\diamond),(T,\diamond))
    \end{aligned}}{%
    \FormulaRefAuto{MonoidSquareInflationProperties}[thm]{5,\allowbreak 2,\allowbreak 10}}
  \proofsteppaired[1,2,3,4,9]{\forall g\in G\,\exists f_g\colon C_g\bij D_{\varphi(g)}}{%
    \FormulaRefAuto{FiniteGroupSquareFiberBijections}[thm]{1,\allowbreak 2,\allowbreak 3,\allowbreak 9,\allowbreak 4}}
  \proofsteppaired[9]{\begin{aligned}[t]
      &\GroupHom{\varphi}{G,\star,e}{H,\diamond,u}\\[-2pt]
      &\land\varphi\colon G\bij H
    \end{aligned}}{%
    \FormulaRefAuto{GroupIsoDef}[def]{9}}
  \proofsteppaired[9]{\SemigroupHom{\varphi}{G,\star}{H,\diamond}}{%
    \FormulaRefAuto{GroupHomDef}[def]{\rAEa{13}}}
  \proofsteppaired[9]{\SemigroupIso{\varphi}{G,\star}{H,\diamond}}{%
    \FormulaRefAuto{SemigroupIsoDef}[def]{14,\allowbreak \rAEb{13}}}
  \proofsteppaired[1,2,3,4,9]{\exists f\,\SemigroupIso{f}{S,\star}{T,\diamond}}{%
    \FormulaRefAuto{InflationFiberIsomorphism}[thm]{11,\allowbreak 15,\allowbreak 12}}
  \proofsteppaired[1,2,3,4,7]{\exists f\,\SemigroupIso{f}{S,\star}{T,\diamond}}{\rEE{8,\allowbreak 9,\allowbreak 16}}
  \proofsteppaired[1,2,3,4]{\exists f\,\SemigroupIso{f}{S,\star}{T,\diamond}}{\rEE{6,\allowbreak 7,\allowbreak 17}}
\end{tabproofpaired}

\paragraph{Die Nichtleerheit des Zielträgers folgt aus den Voraussetzungen.}
Für diese Folgerung wird das Produktquadrat nur auf der bereits als
nichtleer vorausgesetzten Quellhalbgruppe verwendet. Aus einem
Potenzisomorphismus folgt dann auch die Nichtleerheit des Zielträgers.

\par\noindent\begin{minipage}{\linewidth}
\FormulaThmDeltaK[Rekonstruktion ohne Nichtleerheitsannahme am Ziel]{%
  \begin{aligned}[t]
    &\FiniteSemigroupAlg{S}{\star}\dsep\FiniteSemigroupAlg{T}{\diamond}
      \dsep S\neq\varnothing\dsep\GroupAlg{S^2}{\star}{e}\dsep{}\\[-2pt]
    &\SemigroupIso{F}{\PowNE S,\star_{\mathcal P}}{\PowNE T,\diamond_{\mathcal P}}
      \vdash\exists f\,\SemigroupIso{f}{S,\star}{T,\diamond}
  \end{aligned}
}{FiniteGroupSquarePowerSemigroupReconstructionOneSidedNonempty}{%
  \DeltaRow{Mengen}{S\dsep T\dsep e\dsep X\dsep y}
  \DeltaRow{Funktionen}{\star\dsep\diamond\dsep F\dsep f}
}
\end{minipage}\par
\begin{tabproofpaired}
  \proofsteppaired[1]{\FiniteSemigroupAlg{S}{\star}\land\FiniteSemigroupAlg{T}{\diamond}}{\rA}
  \proofsteppaired[2]{S\neq\varnothing}{\rA}
  \proofsteppaired[3]{\GroupAlg{S^2}{\star}{e}}{\rA}
  \proofsteppaired[4]{\SemigroupIso{F}{Q_S,\star_{\mathcal P}}{Q_T,\diamond_{\mathcal P}}}{\rA}
  \proofsteppaired[2]{\PowNE S\neq\varnothing}{\FormulaRefAuto{NonemptyPowerCarrierNonempty}[thm]{2}}
  \proofsteppaired[2]{\exists X\,(X\in\PowNE S)}{%
    \FormulaRefAuto{A\neq\varnothing\eqvdash\exists x\,(x\in A)}[thm]{5}}
  \proofsteppaired[7]{X\in\PowNE S}{\rA}
  \proofsteppaired[4]{F\colon\PowNE S\to\PowNE T}{%
    \FormulaRefAuto{Bijektive Funktion}[def]{\FormulaRefAuto{SemigroupIsoBijectiveAxiom}[ax]{4}}}
  \proofsteppaired[4,7]{F(X)\in\PowNE T}{%
    \FormulaRefAuto{F\colon A\to B\dsep x\in A\vdash F(x)\in B}[thm]{8,\allowbreak 7}}
  \proofsteppaired[4,7]{F(X)\subseteq T\land F(X)\neq\varnothing}{%
    \FormulaRefAuto{NonemptyPowersetMembership}[thm]{9}}
  \proofsteppaired[4,7]{\exists y\,(y\in F(X))}{%
    \FormulaRefAuto{A\neq\varnothing\eqvdash\exists x\,(x\in A)}[thm]{\rAEb{10}}}
  \proofsteppaired[12]{y\in F(X)}{\rA}
  \proofsteppaired[4,7,12]{y\in T}{%
    \FormulaRefAuto{A\subseteq B\dsep x\in A\vdash x\in B}[thm]{\rAEa{10},\allowbreak 12}}
  \proofsteppaired[4,7,12]{T\neq\varnothing}{%
    \FormulaRefAuto{A\neq\varnothing\eqvdash\exists x\,(x\in A)}[thm]{\rEI{13}}}
  \proofsteppaired[4,7]{T\neq\varnothing}{\rEE{11,\allowbreak 12,\allowbreak 14}}
  \proofsteppaired[2,4]{T\neq\varnothing}{\rEE{6,\allowbreak 7,\allowbreak 15}}
  \proofsteppaired[1,2,3,4]{\exists f\,\SemigroupIso{f}{S,\star}{T,\diamond}}{%
    \FormulaRefAuto{FiniteGroupSquarePowerSemigroupReconstruction}[thm]{1,\allowbreak 2,\allowbreak 16,\allowbreak 3,\allowbreak 4}}
\end{tabproofpaired}

Der Satz verlangt keine Erhaltung sämtlicher Einermengen durch \(F\).
Erkennbar sind die Einermengen des Gruppenkerns; außerhalb dieses Kerns
wird eine passende Grundbijektion erst aus den Fasern zusammengesetzt.
Insbesondere umfasst das Ergebnis Nullhalbgruppen: Deren Produktquadrat
\(\{0\}\) ist die triviale Gruppe. Der allgemeine endliche Fall der
Potenzhalbgruppenvermutung bleibt davon getrennt.

\end{document}
